\chapter{Metodología}
La metodología propuesta describe el diseño e implementación de un Autómata Celular (AC) unidimensional para la encriptación y desencriptación de imágenes. El proceso se centra en la aplicación de una variante reversible de la Regla 90, diseñada para manipular los valores de los píxeles de una imagen de forma segura y recuperable.

\section{Representación de la Imagen}
Para aplicar el autómata celular a una imagen, es necesario establecer una correspondencia entre los píxeles de la imagen y las celdas del AC.

\subsection{Objetivo}
Transformar la imagen bidimensional en una representación compatible con un autómata celular.

\subsection{Implementación}
\begin{itemize}
\item Conversión de Imagen: Las imágenes de entrada (e.g., JPEG, PNG) se cargan y, si es necesario, se convierten a un formato de escala de grises para simplificar el manejo de los valores de los píxeles (0-255). Para imágenes a color, se puede optar por encriptar cada canal (R, G, B) de forma independiente o convertir a un espacio de color diferente (e.g., YCbCr) y encriptar el canal de luminancia.
\item Transformación de Píxeles a Bits: La imagen en escala de grises se convierte en una matriz de bits donde cada píxel, representado originalmente por un valor entero de 8 bits (0 a 255), se transforma en su representación binaria fila por fila, expandiendo cada valor en sus 8 bits equivalentes. El resultado es una matriz binaria en la que cada fila de la imagen original se convierte en una secuencia de bits, lo que permite una manipulación más fina de la información a nivel binario para su uso en el autómata celular.
\end{itemize}

\subsection{Ventajas}
\begin{itemize}
\item Permite aplicar un autómata celular unidimensional directamente sobre imágenes bidimensionales al representar cada fila como una secuencia binaria continua.
\item Facilita la manipulación de los datos a nivel de bit, lo cual es útil para operaciones de encriptación más detalladas y controladas.
\item Aprovecha la estructura binaria del autómata celular, alineándola naturalmente con la representación digital de los píxeles (8 bits por píxel).
\item Mantiene la posibilidad de reconstrucción exacta de la imagen original a partir de su representación binaria.
\end{itemize}

\section{Definición de la Regla 90 Reversible}
La Regla 90 es una regla de autómata celular elemental (ECA) conocida por generar patrones complejos, a menudo fractales. Para la encriptación, es crucial que la transformación sea reversible.

\subsection*{Mecanismo de la Regla 90}
Matemáticamente, la Regla 90 se puede expresar como \( C_i(t+1) = C_{i-1} \oplus C_{i+1} \), donde \( \oplus \) denota la operación XOR y los valores son binarios. Pero el uso de esta regla tiene un inconveniente, ya que esta no es reversible por lo que se realiza un pequeño cambio para que sea reversible considerando la fila anterior quedando de la siguiente manera:\( C_i(t+1) = C_{i-1} \oplus C_{i+1} \oplus bits_{i-1}\) 

\subsection*{Importancia de la Reversibilidad}
La reversibilidad es un requisito fundamental para la encriptación de imágenes. Sin ella, la imagen original no podría ser recuperada por el destinatario autorizado, lo que anularía el propósito de la encriptación. Garantiza que la transformación realizada por el AC no pierde información.

\section{Proceso de Encriptación de Imágenes}
La encriptación se logra aplicando la evolución del autómata celular sobre los píxeles de la imagen.

\subsection{Flujo del Proceso}
\begin{itemize}
\item Paso 1: Carga y Preparación de la Imagen: La imagen original se carga y se realiza la transformación de pixeles a bits.
\item Paso 2: Evolución del Autómata: El array de píxeles se considera el estado inicial del autómata. Se aplica la Regla 90 reversible durante un número predefinido de iteraciones (generaciones). En cada iteración, el estado de cada celda (píxel) se actualiza simultáneamente en función de sus vecinos, según la regla.
\item Paso 3: Reconstrucción de la Imagen Encriptada: Una vez completadas las iteraciones, el estado final del autómata (el array de píxeles modificado) se reestructura en la forma bidimensional original de la imagen. Esta es la imagen encriptada.
\end{itemize}

\section{Proceso de Desencriptación de Imágenes}
Dado que la Regla 90 utilizada es reversible, la desencriptación es el proceso inverso de la encriptación.

\subsection{Flujo del Proceso}
\begin{itemize}
\item Paso 1: Carga y Preparación de la Imagen Encriptada: La imagen encriptada se carga y se realiza la transformación de pixeles a bits.
\item Paso 2: Inversión de la Evolución del Autómata: Utilizando la propiedad de reversibilidad de la Regla 90, se invierte el proceso de evolución. Esto implica aplicar la función inversa de la regla de transición un número igual de veces que en la encriptación.
\item Paso 3: Reconstrucción de la Imagen Original: El estado final del autómata después de la inversión se reestructura en la imagen bidimensional original, revelando la imagen desencriptada.
\end{itemize}
